% GENERATED FILE. Edit canonical lessons, then run npm run generate:books.
% canonical-source-hash: fnv1a64:e67557f334b6f95e
% canonical-lessons: AR-C340-hakim, AR-C340-sultan, AR-C340-qabila, AR-C340-balat, AR-C340-tawalla, AR-R340-review-333-340

\chapter{Rulers and tribes}
\label{ch:ar-rulers-and-tribes}
\hlchaptermodality{340}

\begin{chapteropening}
\textbf{By the end of this chapter:} \emph{I can say a ruler, a sultan, a tribe, a royal court and to take office.}

Complete the last lesson of chapter 340: I can say a ruler, a sultan, a tribe, a royal court and to take office.
\end{chapteropening}

\section[\texorpdfstring{\ar{حاكم} (ḥākim)}{حاكم (ḥākim)}]{\ar{حاكم} (ḥākim) — a governor; a ruler}

\label{lesson:AR-C340-hakim}



\begin{quote}
Before the new one: say the Arabic for to lead, then the Arabic for to withdraw.
\end{quote}

\subsection*{You'll want to know: \ar{حاكم}}
\textbf{\ar{حاكم}} — \emph{ḥākim} — \textquotedblleft{}a governor; a ruler\textquotedblright{}.

\begin{quote}
\ar{كان} \ar{الحاكم} \ar{عادلا} \ar{مع} \ar{الناس}.
\end{quote}

\begin{quote}
\emph{kāna al-ḥākim ʿādilan maʿa an-nās.}
\end{quote}

\begin{quote}
The ruler was fair to the people.
\end{quote}

Plural: \ar{حكام}, \emph{ḥukkām}.

\begin{cousinweb}
\textbf{\ar{حاكم}} (\emph{ḥākim}) is built on the native Arabic root \emph{ḥ-k-m} (Lane), \textquotedblleft{}to judge, to rule\textquotedblright{}. \emph{Ḥākim} is the active participle, \textquotedblleft{}one who rules\textquotedblright{}.
\end{cousinweb}

\begin{grammarlens}[title={What you've built}]
One more word for telling what happened.
\end{grammarlens}

\subsection*{Your turn}
\begin{itemize}
\raggedright
  \item \emph{Say it:} \emph{ḥākim}
  \item \emph{Say it:} \emph{ḥākim}, once more
  \item \emph{Say it:} \emph{kāna al-ḥākim ʿādilan maʿa an-nās.} — The ruler was fair to the people.
  \item \emph{Recall:} say the Arabic for to lead, then the Arabic for to withdraw, then say \emph{ḥākim} again
\end{itemize}

\begin{tcolorbox}[breakable,colback=teal!4,colframe=teal!35!black,title={Before you move on}]
What does \ar{حاكم} mean? (\textquotedblleft{}A governor; a ruler\textquotedblright{}.) Say it once more.
\end{tcolorbox}
\section[\texorpdfstring{\ar{سلطان} (sulṭān)}{سلطان (sulṭān)}]{\ar{سلطان} (sulṭān) — a sultan; power}

\label{lesson:AR-C340-sultan}



\begin{quote}
Before the new one: say the Arabic for to withdraw, then the Arabic for a governor.
\end{quote}

\subsection*{You'll want to know: \ar{سلطان}}
\textbf{\ar{سلطان}} — \emph{sulṭān} — \textquotedblleft{}a sultan; power\textquotedblright{}.

\begin{quote}
\ar{سكن} \ar{السلطان} \ar{في} \ar{قصر} \ar{على} \ar{النهر}.
\end{quote}

\begin{quote}
\emph{sakana as-sulṭān fī qaṣr ʿalā an-nahr.}
\end{quote}

\begin{quote}
The sultan lived in a palace on the river.
\end{quote}

Plural: \ar{سلاطين}, \emph{salāṭīn}.

\begin{cousinweb}
\textbf{\ar{سلطان}} (\emph{sulṭān}) is built on the native Arabic root \emph{s-l-ṭ} (Lane), \textquotedblleft{}to have power over\textquotedblright{}. Lane gives \emph{sulṭān} first as power and proof, then the one who holds power; Aramaic \emph{shulṭānā}, power, is usually cited beside it.
\end{cousinweb}

\begin{grammarlens}[title={What you've built}]
One more word for telling what happened.
\end{grammarlens}

\subsection*{Your turn}
\begin{itemize}
\raggedright
  \item \emph{Say it:} \emph{sulṭān}
  \item \emph{Say it:} \emph{sulṭān}, once more
  \item \emph{Say it:} \emph{sakana as-sulṭān fī qaṣr ʿalā an-nahr.} — The sultan lived in a palace on the river.
  \item \emph{Recall:} say the Arabic for to withdraw, then the Arabic for a governor, then say \emph{sulṭān} again
\end{itemize}

\begin{tcolorbox}[breakable,colback=teal!4,colframe=teal!35!black,title={Before you move on}]
What does \ar{سلطان} mean? (\textquotedblleft{}A sultan; power\textquotedblright{}.) Say it once more.
\end{tcolorbox}
\section[\texorpdfstring{\ar{قبيلة} (qabīla)}{قبيلة (qabīla)}]{\ar{قبيلة} (qabīla) — a tribe}

\label{lesson:AR-C340-qabila}



\begin{quote}
Before the new one: say the Arabic for a governor, then the Arabic for a sultan.
\end{quote}

\subsection*{You'll want to know: \ar{قبيلة}}
\textbf{\ar{قبيلة}} — \emph{qabīla} — \textquotedblleft{}a tribe\textquotedblright{}.

\begin{quote}
\ar{عاشت} \ar{القبيلة} \ar{في} \ar{الصحراء}.
\end{quote}

\begin{quote}
\emph{ʿāshat al-qabīla fī aṣ-ṣaḥrāʾ.}
\end{quote}

\begin{quote}
The tribe lived in the desert.
\end{quote}

\begin{cousinweb}
\textbf{\ar{قبيلة}} (\emph{qabīla}) is built on the native Arabic root \emph{q-b-l} (Lane), \textquotedblleft{}to face, to be in front\textquotedblright{}. Lane gives \emph{qabīla} as a company of people from one father.
\end{cousinweb}

\begin{grammarlens}[title={What you've built}]
One more word for telling what happened.
\end{grammarlens}

\subsection*{Your turn}
\begin{itemize}
\raggedright
  \item \emph{Say it:} \emph{qabīla}
  \item \emph{Say it:} \emph{qabīla}, once more
  \item \emph{Say it:} \emph{ʿāshat al-qabīla fī aṣ-ṣaḥrāʾ.} — The tribe lived in the desert.
  \item \emph{Recall:} say the Arabic for a governor, then the Arabic for a sultan, then say \emph{qabīla} again
\end{itemize}

\begin{tcolorbox}[breakable,colback=teal!4,colframe=teal!35!black,title={Before you move on}]
What does \ar{قبيلة} mean? (\textquotedblleft{}A tribe\textquotedblright{}.) Say it once more.
\end{tcolorbox}
\section[\texorpdfstring{\ar{بلاط} (balāṭ)}{بلاط (balāṭ)}]{\ar{بلاط} (balāṭ) — a royal court; paving stones}

\label{lesson:AR-C340-balat}



\begin{quote}
Before the new one: say the Arabic for a sultan, then the Arabic for a tribe.
\end{quote}

\subsection*{You'll want to know: \ar{بلاط}}
\textbf{\ar{بلاط}} — \emph{balāṭ} — \textquotedblleft{}a royal court; paving stones\textquotedblright{}.

\begin{quote}
\ar{كان} \ar{الشاعر} \ar{يعمل} \ar{في} \ar{بلاط} \ar{الملك}.
\end{quote}

\begin{quote}
\emph{kāna ash-shāʿir yaʿmalu fī balāṭ al-malik.}
\end{quote}

\begin{quote}
The poet worked at the king’s court.
\end{quote}

At home, \ar{بلاط} is floor tiles.

\begin{cousinweb}
\textbf{\ar{بلاط}} (\emph{balāṭ}) is given by Lane under \emph{b-l-ṭ} as stones laid for paving. It was borrowed from Aramaic \emph{pallāṭīn} (EtymDB); the sense of a royal court is usually traced further to Latin \emph{palatium}, a palace.
\end{cousinweb}

\begin{grammarlens}[title={What you've built}]
One more word for telling what happened.
\end{grammarlens}

\subsection*{Your turn}
\begin{itemize}
\raggedright
  \item \emph{Say it:} \emph{balāṭ}
  \item \emph{Say it:} \emph{balāṭ}, once more
  \item \emph{Say it:} \emph{kāna ash-shāʿir yaʿmalu fī balāṭ al-malik.} — The poet worked at the king’s court.
  \item \emph{Recall:} say the Arabic for a sultan, then the Arabic for a tribe, then say \emph{balāṭ} again
\end{itemize}

\begin{tcolorbox}[breakable,colback=teal!4,colframe=teal!35!black,title={Before you move on}]
What does \ar{بلاط} mean? (\textquotedblleft{}A royal court; paving stones\textquotedblright{}.) Say it once more.
\end{tcolorbox}
\section[\texorpdfstring{\ar{تولى} (tawallā)}{تولى (tawallā)}]{\ar{تولى} (tawallā) — to take office; to take charge}

\label{lesson:AR-C340-tawalla}



\begin{quote}
Before the new one: say the Arabic for a tribe, then the Arabic for a royal court.
\end{quote}

\subsection*{You'll want to know: \ar{تولى}}
\textbf{\ar{تولى}} — \emph{tawallā} — \textquotedblleft{}to take office; to take charge\textquotedblright{}.

\begin{quote}
\ar{تولى} \ar{الأمير} \ar{الحكم} \ar{بعد} \ar{أبيه}.
\end{quote}

\begin{quote}
\emph{tawallā al-amīr al-ḥukm baʿda abīhi.}
\end{quote}

\begin{quote}
The prince took power after his father.
\end{quote}

Present: \ar{يتولى}, \emph{yatawallā}.

\ar{تولى} \ar{الأمر}: he took charge of the matter.

\begin{cousinweb}
\textbf{\ar{تولى}} (\emph{tawallā}) is built on the native Arabic root \emph{w-l-y} (Lane), \textquotedblleft{}to be near; to have charge of\textquotedblright{}, the root of \emph{walī}, a guardian. It is Form V, \textquotedblleft{}to take charge\textquotedblright{}.
\end{cousinweb}

\begin{grammarlens}[title={What you've built}]
One more word for telling what happened.
\end{grammarlens}

\subsection*{Your turn}
\begin{itemize}
\raggedright
  \item \emph{Say it:} \emph{tawallā}
  \item \emph{Say it:} \emph{tawallā}, once more
  \item \emph{Say it:} \emph{tawallā al-amīr al-ḥukm baʿda abīhi.} — The prince took power after his father.
  \item \emph{Recall:} say the Arabic for a tribe, then the Arabic for a royal court, then say \emph{tawallā} again
\end{itemize}

\begin{tcolorbox}[breakable,colback=teal!4,colframe=teal!35!black,title={Before you move on}]
What does \ar{تولى} mean? (\textquotedblleft{}To take office; to take charge\textquotedblright{}.) Say it once more.
\end{tcolorbox}
\section[(dialogue)]{Words from chapters 333–340 — a review}

\label{lesson:AR-R340-review-333-340}



\begin{quote}
Say the Arabic for to go under and a sailor.
\end{quote}

\subsection*{Your turn}
Say these aloud:

\begin{itemize}
\raggedright
  \item to go under, a sailor, a fleet, a pirate, to capture
  \item a trial, a hearing, the defence, to convict, innocence
  \item a magician, an ogre, a dragon, a giant, a dwarf
  \item to gift, to receive guests, to bless, a dowry, the eve
  \item a will, an inheritance, to inherit, an heir, a widow
  \item a eulogy, a tomb, a death, a cemetery, a coffin
  \item to invade, to attack, to occupy, to lead, to withdraw
  \item a governor, a sultan, a tribe, a royal court, to take office
\end{itemize}

\begin{tcolorbox}[breakable,colback=teal!4,colframe=teal!35!black,title={Before you move on}]
To go under? (\textbf{\ar{غرق}}, \emph{ghariqa}.) A sailor? (\textbf{\ar{بحار}}, \emph{baḥḥār}.) A fleet? (\textbf{\ar{أسطول}}, \emph{usṭūl}.) A pirate? (\textbf{\ar{قرصان}}, \emph{qurṣān}.) To capture? (\textbf{\ar{أسر}}, \emph{asara}.) A trial? (\textbf{\ar{محاكمة}}, \emph{muḥākama}.) A hearing? (\textbf{\ar{جلسة}}, \emph{jalsa}.) The defence? (\textbf{\ar{دفاع}}, \emph{difāʿ}.) To convict? (\textbf{\ar{أدان}}, \emph{adāna}.) Innocence? (\textbf{\ar{براءة}}, \emph{barāʾa}.) A magician? (\textbf{\ar{ساحر}}, \emph{sāḥir}.) An ogre? (\textbf{\ar{غول}}, \emph{ghūl}.) A dragon? (\textbf{\ar{تنين}}, \emph{tinnīn}.) A giant? (\textbf{\ar{عملاق}}, \emph{ʿimlāq}.) A dwarf? (\textbf{\ar{قزم}}, \emph{qazam}.) To gift? (\textbf{\ar{أهدى}}, \emph{ahdā}.) To receive guests? (\textbf{\ar{استقبل}}, \emph{istaqbala}.) To bless? (\textbf{\ar{بارك}}, \emph{bāraka}.) A dowry? (\textbf{\ar{مهر}}, \emph{mahr}.) The eve? (\textbf{\ar{عشية}}, \emph{ʿashiyya}.) A will? (\textbf{\ar{وصية}}, \emph{waṣiyya}.) An inheritance? (\textbf{\ar{ميراث}}, \emph{mīrāth}.) To inherit? (\textbf{\ar{ورث}}, \emph{waritha}.) An heir? (\textbf{\ar{وارث}}, \emph{wārith}.) A widow? (\textbf{\ar{أرملة}}, \emph{armala}.) A eulogy? (\textbf{\ar{تأبين}}, \emph{taʾbīn}.) A tomb? (\textbf{\ar{ضريح}}, \emph{ḍarīḥ}.) A death? (\textbf{\ar{وفاة}}, \emph{wafāh}.) A cemetery? (\textbf{\ar{مقبرة}}, \emph{maqbara}.) A coffin? (\textbf{\ar{نعش}}, \emph{naʿsh}.) To invade? (\textbf{\ar{غزا}}, \emph{ghazā}.) To attack? (\textbf{\ar{هاجم}}, \emph{hājama}.) To occupy? (\textbf{\ar{احتل}}, \emph{iḥtalla}.) To lead? (\textbf{\ar{قاد}}, \emph{qāda}.) To withdraw? (\textbf{\ar{انسحب}}, \emph{insaḥaba}.) A governor? (\textbf{\ar{حاكم}}, \emph{ḥākim}.) A sultan? (\textbf{\ar{سلطان}}, \emph{sulṭān}.) A tribe? (\textbf{\ar{قبيلة}}, \emph{qabīla}.) A royal court? (\textbf{\ar{بلاط}}, \emph{balāṭ}.) To take office? (\textbf{\ar{تولى}}, \emph{tawallā}.)
\end{tcolorbox}
